% Осень26!A22; нумерация внутри столбца без пустых строк.
% Источник: лекция 2, §14 (строка 348)
% Источник: лекция 2, §15 (строка 384)
\ticket{19}{Теорема об эквивалентности счетной аддитивности и счетной полуаддитивности}
\label{ticket:19}

\begin{definition}
Объём $\mu$ на полукольце $\mathcal P$ называется \textbf{счётно полуаддитивным}, если для любых $A,A_i\in\mathcal P$
\[
  A\subset\bigcup_{i=1}^{\infty}A_i
  \quad\Longrightarrow\quad
  \mu(A)\le\sum_{i=1}^{\infty}\mu(A_i).
\]
\end{definition}

\begin{theorem}
Объём на полукольце счётно аддитивен тогда и только тогда, когда он счётно полуаддитивен.
\end{theorem}

\begin{proof}
Пусть $\mu$ счётно аддитивна и $A\subset\bigcup_i A_i$.
Как в \hyperref[ticket:18]{доказательстве конечной полуаддитивности}, положим
\[
  C_k=(A\cap A_k)\setminus\bigcup_{i<k}(A\cap A_i)
      =\bigsqcup_j D_{kj},\qquad D_{kj}\in\mathcal P.
\]
При каждом $k$ разбиение конечно, все $D_{kj}$ попарно не пересекаются и $A=\bigsqcup_{k,j}D_{kj}$.
По счётной аддитивности и усиленной монотонности
\[
  \mu(A)=\sum_k\sum_j\mu(D_{kj})\le\sum_k\mu(A_k).
\]

Обратно, пусть выполнена счётная полуаддитивность и $A=\bigsqcup_{i=1}^{\infty}A_i\in\mathcal P$.
Она даёт $\mu(A)\le\sum_i\mu(A_i)$.
При каждом $n$ по усиленной монотонности $\sum_{i=1}^n\mu(A_i)\le\mu(A)$.
Переходя к пределу, получаем противоположное неравенство и требуемое равенство.
\end{proof}
